\section{related work}
\label{sec:relatedwork}
\subsection{View Selection and Next Best View}

View selection and \ac{NBV} have received significant attention as approaches for identifying the most informative next viewpoint and selecting an optimal set of input images that effectively summarize the scene content. Previous approaches in radiance field-based \ac{NBV} have included methods that treat uncertainty directly as a learnable parameter \cite{feng2024naruto} or exploit heuristic information gain functions \cite{chen2025activegamer}. However, these methods fail to fully capture epistemic uncertainty arising from insufficient knowledge. 
% \cite{klasson2024sources}

Some approaches \cite{goli2024bayes, jiang2024fisherrf, wilson2025pop, strong2024next}, on the other hand, have proposed quantifying uncertainty using the Hessian matrix of the negative log-likelihood and successfully applied to both \ac{NeRF} and \ac{3DGS}. Specifically, FisherRF \cite{jiang2024fisherrf} directly measured uncertainty by leveraging the Jacobian of rendering outputs with respect to the Gaussian primitives in \ac{3DGS}. Subsequently, POp-GS \cite{wilson2025pop} introduced concepts from optimal experimental design \cite{kiefer1974general} into information gain calculations, partially accounting for parameter correlations. Next Best Sense \cite{strong2024next} further expanded its applicability by incorporating depth information into the information gain.

% 기존 연구에서는 radiance field에서 next best view를 수행하기 위해 uncertainty를 직접 learnable parameter로 설정하는 접근법도 제안되었지만, 이 방식은 지식의 부족으로 인한 epistemic uncertainty를 제대로 반영하지 못하는 한계가 있다. 반면, 최근에는 negative log likelihood의 2차 미분인 Hessian matrix를 이용하여 uncertainty를 정량적으로 분석하는 연구가 제안되었고, 이는 NeRF와 3DGS 모두에 성공적으로 적용되었다. 특히, FisherRF는 3DGS의 Gaussian primitive에 대한 rendering 결과의 Jacobian을 이용해 직접 uncertainty를 측정했다. 이후 Pop-GS는 optimal experimental design 개념을 활용하여 information gain 계산에 parameter 간의 correlation을 부분적으로 반영했으며, Next Best Sense는 depth 정보까지 information gain 계산에 포함하여 활용 가능성을 더욱 확장했다.

However, conventional information-gain-based \ac{NBV} methods often display a strong bias toward regions that have already been sufficiently observed, while leaving underexplored areas inadequately covered. This imbalance between exploration and exploitation limits their effectiveness in completing the scene understanding. We address this by incorporating per-Gaussian confidence weighting to guide viewpoint selection more evenly.

% However, conventional information-gain-based \ac{NBV} methods tend to overly focus on regions already sufficiently observed, neglecting areas that genuinely require further exploration. To mitigate this issue, our method employs per-Gaussian confidence weights derived from each Gaussian's one-hot object vector, ensuring that regions with limited observations receive higher priority. This approach effectively balances exploration and exploitation during viewpoint selection.

% 그러나 기존의 information gain 기반 NBV 방법은 이미 충분히 관찰되어 더 이상 추가적인 관찰이 필요하지 않은 영역에 계속 집중하는 경향이 있다. 이는 잘 관찰된 Gaussians들이 렌더링 품질에 크게 기여하기 때문에 information gain이 높게 평가되어, 정작 occluded region에 대한 exploration이 저해되는 문제를 야기한다. 이러한 문제를 해결하기 위해 본 연구에서는 Gaussian별로 관찰의 confidence를 명시적으로 고려하여 confidence-aware information gain을 제안한다. 이는 기존 NBV 방식과 달리 Gaussian의 object assignment 정보에서 유도된 confidence weight를 이용하여, 관찰이 부족한 영역의 Gaussians에 더 높은 가중치를 부여하고 exploration과 exploitation의 균형을 효과적으로 유지할 수 있도록 한다.

\subsection{Feature Embedded 3D Representation for Manipulation}

% Owing to its explicit representation and ease of customization, \ac{3DGS} has emerged as a powerful framework for comprehensive 3D scene understanding, spanning from low-level geometric reconstruction to high-level object localization and identification.

Recent studies \cite{kobayashi2022decomposing, kerr2023lerf} have demonstrated that applying feature distillation to volumetric rendering can jointly capture semantic and geometric information, making it highly suitable for robotic manipulation. This idea has also been extended to \ac{3DGS}, where several works \cite{qiu2024feature, qin2024langsplat} have customized the CUDA rasterizer to incorporate features from various language-based models such as CLIP \cite{radford2021learning}, SAM \cite{kirillov2023segment}, and GroundingDINO \cite{liu2024grounding}. GraspSplats \cite{ji2024graspsplats} leverages masks obtained from MobileSAM \cite{zhang2023faster} and their CLIP embeddings to design a hierarchical feature representation, enabling zero-shot object manipulation. However, directly injecting such features into Gaussians is challenging because their high dimensionality complicates optimization, slows processing, and makes it difficult to distinguish among semantically similar instances.
% However, directly injecting such features into Gaussians poses significant challenges, as the high dimensionality of the features complicates optimization, slows down overall processing, and makes it difficult to distinguish between instances sharing the same semantics.

% DFF와 LERF의 insight에서 출발하여 최근 연구자들은 feature distillation을 적용시킨 volumetric rendering이 semantic 정보와 geometry 정보를 함께 주기 때문에 robotic manipulation에 매우 적합하다는 것을 알아내었다. 이는 3DGS에도 적용되어 여러 연구에서 cuda rasterizer를 custom하여 clip, sam, groundingdino 등 여러 language-based model들의 feature를 넣어 이를 구현하였다. GraspSplats에서는 mobilesam으로 얻은 mask와 그들의 clip embedding으로 hierarchical feature를 설계하여 zero-shot object manipulation을 선보였다. 그러나, gaussian 들에 직접 이 feature를 inject하는 것은 feature의 dimension이 커서 optimzie가 어려울 뿐만 아니라 전반적인 속도를 늦추고, 같은 semantic의 instance들을 구분하기 어렵게 한다.

Accordingly, research has also emerged focusing on faster reconstruction using simpler features. Gaussian Grouping \cite{ye2024gaussian} introduced identity encodings to cluster Gaussians belonging to the same object, while ObjectGS \cite{zhu2025objectgs} and TRAN-D \cite{kim2025tran} segmented and tracked objects using Grounded SAM \cite{ren2024grounded} and then optimized a one-hot vector for each object. Similarly, we employ a one-hot vector as supervision for objects, and we further extend its use by leveraging it as a confidence measure integrated into NBV planning. This also serves as the cornerstone of our object-centric reconstruction.

% 이에 따라 좀 더 간단한 feature로 빠르게 reconstruction을 하기 위한 연구들도 나왔다. Gaussian Grouping에서는 introduced identity encodings to cluster Gaussians belonging to the same object, and objectgs와 TRAN-D는 grounded sam으로 물체들을 segment 및 tracking 한 뒤 object별의 one hot vector를 optimize했다. 우리 또한 object의 one hot vector를 supervision으로 사용하지만, 우리는 거기에서 더 나아가 이를 confidence로 활용하여 NBV planning과 결합한다.
